% !TeX root = ../概率论与数理统计笔记.tex
% 由主文档通过 \input 载入；请编译主文档。
\setcounter{chapter}{13} % 保留原书章号。
\chapter{假设检验的原理与方法}
\label{chap:ps-testing}

本章研究如何在控制错误概率的前提下，根据数据评价关于总体的假设。
先修内容为抽样分布、似然和区间估计；学习目标是完整陈述检验问题、选择统计量、校准拒绝域并解释结论。

\section{假设、错误概率与功效}
\label{sec:ps-testing-framework}
定义\emph{原假设}（null hypothesis）$H_0:\theta\in\Theta_0$ 与\emph{备择假设}（alternative hypothesis）$H_1:\theta\in\Theta_1$，其中 $\Theta_0,\Theta_1$ 为不相交的参数子集。
区分简单假设、复合假设、单侧和双侧问题，定义检验统计量、拒绝域、第一类和第二类错误。
水平 $\alpha$ 的检验要求在原假设参数集合上拒绝概率的上确界不超过 $\alpha$；\emph{功效}（power）描述在备择参数处拒绝原假设的概率。

\section{显著性水平与 p 值}
\label{sec:ps-p-values}
解释\emph{p 值}（p-value）作为在指定检验及原假设校准方式下衡量观测结果极端程度的量。
在嵌套拒绝域框架下说明它与可拒绝的最小显著性水平之间的关系，离散和复合原假设情形需注明定义方式。
不把 p 值解释为原假设为真的概率，也不把未拒绝解释为证实原假设。

\section{正态总体与比例的常用检验}
\label{sec:ps-classical-tests}
推导单样本与双样本均值检验、配对检验、正态方差检验及独立正态样本方差比检验。
讨论二项比例的精确检验和大样本检验，明确每种方法的模型条件、统计量分布与自由度。
对同一问题比较检验结论与相应置信区间，建立匹配方法下的对偶关系。

\section{最强检验与似然比检验}
\label{sec:ps-likelihood-ratio-tests}
在简单原假设对简单备择假设的条件下陈述\emph{奈曼--皮尔逊引理}（Neyman--Pearson lemma），讨论离散情形达到指定水平可能需要随机化。
介绍\emph{似然比检验}（likelihood-ratio test），在约束参数空间与完整参数空间分别最大化似然。
渐近卡方结论仅在适当可识别、光滑、内点和信息非退化条件下使用，边界参数或不规则模型另行分析。

\section{检验设计、效应量与多重比较}
\label{sec:ps-power-multiplicity}
研究样本量、显著性水平、效应大小与功效之间的关系，区分统计显著与问题中的实际差异大小。
说明预先确定分析方案、报告区间和效应量的作用；介绍多次检验对总体错误概率的影响。
以邦费罗尼方法（Bonferroni method）说明一类基本控制方式，并明确所控制的是族错误率。

\section{例题、习题与学习检查}
规划均值是否超过给定阈值、两独立总体均值差、配对测量差、伯努利模型的最强检验。
习题要求写齐假设、条件、统计量、拒绝规则、观测结果与语境解释，并计算选定备择下的功效。
学习检查强调检验针对模型中的假设，数据质量与模型适用性不能由一个 p 值代替。






